% =================================================
% Теоретический лист. Урок 2
% =================================================

\begin{center}
  {\fontsize{25}{30}\selectfont\bfseries\sffamily\color{Primary}
    Классы. Чтение и запись чисел}

  \vspace{0.3em}

  {\large\sffamily\color{GrayText}
    Урок 2 из 3. Обозначение натуральных чисел}
\end{center}

\vspace{0.45em}

\section{Зачем нужны классы}

У больших чисел много разрядов. Чтобы запись было легче читать, цифры
объединяют в группы по три. Такие группы называют \textbf{классами}.

\splitclassesfigure

\begin{propertybox}[Как разделить число на классы]
  Начинайте с правого края записи и отделяйте по три цифры. Самая левая
  группа может содержать одну, две или три цифры.
\end{propertybox}

Например:
\[
  7304052018=7\,|\,304\,|\,052\,|\,018.
\]
Пробелы между классами не меняют число: они делают его запись нагляднее.

\section{Какие бывают классы}

Классы считают справа налево:
\[
  \underbrace{\text{единицы}}_{\text{I класс}},\qquad
  \underbrace{\text{тысячи}}_{\text{II класс}},\qquad
  \underbrace{\text{миллионы}}_{\text{III класс}},\qquad
  \underbrace{\text{миллиарды}}_{\text{IV класс}}.
\]

В каждом классе три разряда: единицы, десятки и сотни этого класса.

\classtablefigure

Для числа \(7\,304\,052\,018\):
\begin{itemize}
  \item класс миллиардов содержит число \(7\);
  \item класс миллионов --- \(304\);
  \item класс тысяч --- \(52\), но в записи класса стоит \(052\);
  \item класс единиц --- \(18\), но в записи класса стоит \(018\).
\end{itemize}

\begin{notebox}[Почему внутри класса записывают три цифры]
  Все классы, кроме самого левого, занимают ровно три разряда. Поэтому
  \(52\) тысячи записывают как группу \(052\), а \(18\) единиц --- как
  группу \(018\).
\end{notebox}

\section{Как прочитать многозначное число}

\readingroutefigure

\begin{conclusionbox}[Алгоритм чтения]
  \begin{enumerate}
    \item Разделите число справа налево на классы по три цифры.
    \item Начните с крайнего левого ненулевого класса.
    \item Прочитайте число внутри класса и назовите класс.
    \item Если класс содержит \(000\), пропустите его при чтении.
    \item Число класса единиц прочитайте без слов «класс единиц».
  \end{enumerate}
\end{conclusionbox}

Пример:
\[
  7\,304\,052\,018
\]
читают так:

\begin{center}
  \textbf{семь миллиардов триста четыре миллиона\newline
  пятьдесят две тысячи восемнадцать}.
\end{center}

\begin{notebox}[Форма слова «тысяча»]
  Говорят: \textbf{одна тысяча}, \textbf{две тысячи}, но
  \textbf{один миллион}, \textbf{два миллиона}.
\end{notebox}

\section{Нулевой класс}

\zeroclassfigure

\begin{warningbox}[Нулевой класс нельзя удалить из записи]
  Число \(9\,000\,406\) не равно \(9\,406\). В первом числе цифра \(9\)
  обозначает миллионы, во втором --- тысячи.
\end{warningbox}

Ещё один пример:
\[
  2\,005\,008
\]
читают: \textbf{два миллиона пять тысяч восемь}. Нули внутри групп
\(005\) и \(008\) сохраняют отсутствующие разряды.

\section{Как записать число цифрами}

\writingclassesfigure

\begin{conclusionbox}[Алгоритм записи]
  \begin{enumerate}
    \item Определите, какие классы названы словами.
    \item Запишите число каждого класса отдельной группой.
    \item Каждый класс справа от первого дополните слева нулями до трёх цифр.
    \item Если какой-либо промежуточный класс не назван, запишите группу \(000\).
    \item Соедините группы, оставив между ними небольшие пробелы.
  \end{enumerate}
\end{conclusionbox}

Например:
\[
  \text{сорок два миллиона семь тысяч девяносто}
\]
\[
  42\;|\;007\;|\;090=42\,007\,090.
\]

Если класс тысяч не назван:
\[
  \text{триста пять миллионов восемьсот два}
\]
\[
  305\;|\;000\;|\;802=305\,000\,802.
\]

\begin{warningbox}[Типичная ошибка]
  Неверно:
  \[
    \text{шесть миллионов сорок два}=6\,042.
  \]
  Число \(42\) относится к классу единиц, а между классами миллионов и
  единиц должен стоять пустой класс тысяч:
  \[
    \text{шесть миллионов сорок два}=6\,000\,042.
  \]
\end{warningbox}

\section{Как проверить запись}

\begin{trainingtask}[Обратное действие]
  Прочитайте полученное число вслух. Если слова совпали с условием, запись
  выполнена правильно.

  Например:
  \[
    53\,008\,004
  \]
  читается: «пятьдесят три миллиона восемь тысяч четыре».
\end{trainingtask}

\begin{practicebox}[Проверьте себя]
  \begin{enumerate}
    \item Разделите на классы число \(4805070204\).
    \item Какое число находится в классе тысяч числа \(325\,041\,708\)?
    \item Прочитайте число \(6\,040\,019\).
    \item Запишите цифрами: девять миллиардов три миллиона двадцать тысяч четыре.
  \end{enumerate}
\end{practicebox}

\begin{notebox}[Ответы]
  1) \(4\,|\,805\,|\,070\,|\,204\); \quad
  2) \(41\); \quad
  3) шесть миллионов сорок тысяч девятнадцать; \quad
  4) \(9\,003\,020\,004\).
\end{notebox}

\begin{questionsbox}
  \begin{enumerate}
    \item С какого края число делят на классы?
    \item Сколько разрядов содержит каждый класс?
    \item Почему класс \(005\) нельзя записать как \(5\) внутри многозначного числа?
    \item Как поступить, если промежуточный класс не назван?
  \end{enumerate}
\end{questionsbox}
